\section{SPARQL Queries}
\label{sec:sparql}

All queries were executed against the Apache Fuseki \ac{sparql}~1.1 endpoint described in Section~\ref{sec:deployment} (\texttt{http://localhost:3030/cineexplorer/query}). Query files and \ac{csv} result exports are available in the \texttt{sparql/} directory of the repository.

\subsection{Query Overview}

Table~\ref{tab:sparql-overview} maps each query to the \ac{sparql}~1.1 feature it demonstrates.

\begin{table}[H]
\centering
\small
\begin{tabular}{llp{6cm}}
\hline
\textbf{Query} & \textbf{File} & \textbf{SPARQL feature(s)} \\
\hline
Q1  & \texttt{q01\_films\_genres.sparql}          & SELECT, multi-join, ORDER BY \\
Q2  & \texttt{q02\_roles\_count.sparql}           & COUNT DISTINCT, GROUP BY, HAVING \\
Q3  & \texttt{q03\_worked\_together.sparql}       & Subquery, COUNT DISTINCT, ORDER BY, LIMIT \\
Q4  & \texttt{q04\_actors\_multi\_genre.sparql}   & Subquery with HAVING, outer GROUP BY \\
Q5  & \texttt{q05\_no\_director.sparql}           & UNION, FILTER NOT EXISTS (negation) \\
Q6  & \texttt{q06\_bacon\_number.sparql}          & Property paths (\texttt{/}, \texttt{\{1,2\}}) \\
Q7  & \texttt{q07\_wikidata\_federation.sparql}   & SERVICE (federated query), BIND, REPLACE \\
Q8  & \texttt{q08\_series\_seasons.sparql}        & COUNT DISTINCT, OPTIONAL, GROUP BY \\
Q9  & \texttt{q09\_director\_never\_actor.sparql} & MINUS (set subtraction) \\
Q10 & \texttt{q10\_similar\_titles.sparql}        & Subquery, multi-join, COUNT DISTINCT \\
\hline
\end{tabular}
\caption{\ac{sparql} queries and the \ac{sparql}~1.1 features they demonstrate.}
\label{tab:sparql-overview}
\end{table}

\subsection{Q1 — Films with Genres and Runtime}

This query retrieves all films in the knowledge graph together with their genres and runtime, ordered from longest to shortest. It demonstrates multi-property joins across the \texttt{ce:Film}~$\to$~\texttt{ce:Genre} chain and serves as a baseline sanity check for the deployment.

\begin{lstlisting}[language=SPARQL]
PREFIX ce: <http://cineexplorer.local/ontology#>

SELECT ?title ?genre ?runtime
WHERE {
  ?film a ce:Film ;
        ce:workTitle ?title ;
        ce:runtimeMinutes ?runtime ;
        ce:hasGenre ?g .
  ?g ce:genreName ?genre .
}
ORDER BY DESC(?runtime)
\end{lstlisting}

The query returns 11{,}252 rows (each film appears once per genre, since \texttt{ce:hasGenre} is multi-valued). The three longest films are \emph{Gangs of Wasseypur}~(321\,min, Action/Comedy/Crime), \emph{Zack Snyder's Justice League}~(242\,min), and \emph{Gone with the Wind}~(238\,min).

\subsection{Q2 — Persons per Participation Role}

This query counts how many distinct persons hold each participation role and filters to roles with more than five persons. It demonstrates \texttt{COUNT DISTINCT}, \texttt{GROUP BY}, and \texttt{HAVING}.

\begin{lstlisting}[language=SPARQL]
PREFIX ce: <http://cineexplorer.local/ontology#>

SELECT ?role (COUNT(DISTINCT ?person) AS ?personCount)
WHERE {
  ?par a ce:Participation ;
       ce:participationRole ?role ;
       ce:playedBy ?person .
}
GROUP BY ?role
HAVING (COUNT(DISTINCT ?person) > 5)
ORDER BY DESC(?personCount)
\end{lstlisting}

\begin{table}[H]
\centering
\begin{tabular}{lr}
\hline
\textbf{Role} & \textbf{Persons} \\
\hline
actor               & 22{,}705 \\
writer              & 5{,}383  \\
producer            & 4{,}093  \\
director            & 1{,}825  \\
editor              & 1{,}734  \\
composer            & 1{,}308  \\
production\_designer & 1{,}303  \\
cinematographer     & 1{,}261  \\
casting\_director    & 1{,}092  \\
self                & 333      \\
archive\_footage     & 47       \\
\hline
\end{tabular}
\caption{Q2 result: participation roles with more than 5 distinct persons.}
\label{tab:q2-results}
\end{table}

The result (Table~\ref{tab:q2-results}) shows eleven roles above the threshold. Note that \texttt{actor} subsumes both male and female performers in the migrated dataset: at \ac{etl} time the \ac{imdb} categories \texttt{actor} and \texttt{actress} are collapsed into a single \texttt{actor} value, so the 22{,}705 figure here covers what \ac{imdb} originally labels as separate categories.

\subsection{Q3 — Persons Who Worked Together Most}

This query identifies the top 10 pairs of persons who collaborated on the most shared titles. The inner subquery aggregates and ranks all pairs; the outer query resolves their names. \texttt{FILTER(STR(?p1) < STR(?p2))} ensures each pair is counted once rather than both~(A,~B) and~(B,~A).

\begin{lstlisting}[language=SPARQL]
PREFIX ce: <http://cineexplorer.local/ontology#>

SELECT ?name1 ?name2 ?sharedTitles
WHERE {
  {
    SELECT ?p1 ?p2 (COUNT(DISTINCT ?title) AS ?sharedTitles)
    WHERE {
      ?p1 ce:workedFor ?title .
      ?p2 ce:workedFor ?title .
      FILTER(STR(?p1) < STR(?p2))
    }
    GROUP BY ?p1 ?p2
    ORDER BY DESC(?sharedTitles)
    LIMIT 10
  }
  ?p1 ce:personName ?name1 .
  ?p2 ce:personName ?name2 .
}
ORDER BY DESC(?sharedTitles)
\end{lstlisting}

The top-10 pairs are dominated by long-running casting director and producer partnerships, who repeatedly co-credit on the same productions. The first three rows are Deborah Aquila with Tricia Wood (63 shared titles, both casting directors), Janet Hirshenson with Jane Jenkins (55 titles), and Tim Bevan with Eric Fellner (52 titles, the founders of Working Title Films). The migrated \ac{kg} --- where each person averages 2.78 credits across the slice --- supports collaboration counts well above one, which the previous curated sample of 174 titles could not.

\subsection{Q4 — Actors Spanning Three or More Genres}

This query finds actors who appeared in titles covering at least three different genres. A subquery with \texttt{HAVING} first identifies qualifying actors; the outer query counts their genres. Using a subquery avoids computing genre counts for all 38{,}067 persons before filtering.

\begin{lstlisting}[language=SPARQL]
PREFIX ce: <http://cineexplorer.local/ontology#>

SELECT ?name (COUNT(DISTINCT ?genre) AS ?genreCount)
WHERE {
  {
    SELECT ?person
    WHERE {
      ?person a ce:Actor .
      ?person ce:actedIn ?title .
      ?title ce:hasGenre ?g .
      ?g ce:genreName ?genre .
    }
    GROUP BY ?person
    HAVING (COUNT(DISTINCT ?genre) >= 3)
  }
  ?person ce:personName ?name .
  ?person ce:actedIn ?title .
  ?title ce:hasGenre ?g .
  ?g ce:genreName ?genre .
}
GROUP BY ?person ?name
ORDER BY DESC(?genreCount)
\end{lstlisting}

The query returns 17{,}832 actor-genre rows: at the genre-spanning threshold of three, the migrated \ac{kg} yields a wide set of qualifying actors (the row count is one row per actor-genre pair retained for the outer COUNT). The top of the result --- actors spanning the highest number of distinct genres --- is dominated by veteran performers with long \ac{imdb} careers.

\subsection{Q5 — Titles with No Director Credited}

This query identifies films and series that have no director credited in the knowledge graph, using \texttt{FILTER NOT EXISTS} as negation as failure.

\begin{lstlisting}[language=SPARQL]
PREFIX ce: <http://cineexplorer.local/ontology#>

SELECT ?title ?type
WHERE {
  { ?work a ce:Film .   BIND("Film"   AS ?type) }
  UNION
  { ?work a ce:Series . BIND("Series" AS ?type) }
  ?work ce:workTitle ?title .
  FILTER NOT EXISTS {
    ?person ce:directed ?work .
  }
}
ORDER BY ?type ?title
\end{lstlisting}

The query returns 626 titles (predominantly TV series) without a credited director. TV series frequently list directors only at the episode level rather than at the series level, which explains the dominance of \texttt{Series} rows in the output. Under the Open World Assumption, the absence of a \texttt{ce:directed} triple does not mean the work has no director --- only that this information is not represented in the knowledge graph.

\subsection{Q6 — Collaboration Distance via Property Paths}

This query explores the collaboration graph using \ac{sparql}~1.1 property paths. Russell Crowe (\texttt{nm0000128}) is used as the source for two reasons: he is well represented in the migrated \ac{kg} (558 direct collaborators across his credited filmography in the slice), and he is one of the historical anchors of the project's running example through \emph{Les Mis\'{e}rables}~(2012).

\textbf{Direct collaborators (1-hop):}

\begin{lstlisting}[language=SPARQL]
PREFIX ce: <http://cineexplorer.local/ontology#>

SELECT DISTINCT ?collaboratorName
WHERE {
  BIND(<http://cineexplorer.local/data/person/nm0000128> AS ?source)
  ?source ce:workedFor ?title .
  ?collaborator ce:workedFor ?title .
  FILTER(?collaborator != ?source)
  ?collaborator ce:personName ?collaboratorName .
}
ORDER BY ?collaboratorName
\end{lstlisting}

\textbf{Persons reachable within 2~hops:}

\begin{lstlisting}[language=SPARQL]
PREFIX ce: <http://cineexplorer.local/ontology#>

SELECT DISTINCT ?reachableName
WHERE {
  BIND(<http://cineexplorer.local/data/person/nm0000128> AS ?source)
  ?source (ce:workedFor/ce:employed){1,2} ?reachable .
  FILTER(?reachable != ?source)
  ?reachable ce:personName ?reachableName .
}
ORDER BY ?reachableName
\end{lstlisting}

The path \texttt{(ce:workedFor/ce:employed)\{1,2\}} traverses one or two hops, where each hop corresponds to a shared title. The 1-hop query returns 558 direct collaborators of Russell Crowe; the 2-hop query returns 21{,}214 distinct persons reachable within two collaborations. The 38$\times$ growth from 1-hop to 2-hop is the headline structural property of the migrated dataset --- the previous curated sample of 174 titles had no path-multiplying intermediaries because no person was credited on more than one title.

\subsection{Q7 — Wikidata Federation}

This query enriches CineExplorer persons with birth dates from Wikidata via a federated \texttt{SERVICE} call. Wikidata property P345 stores \ac{imdb} person identifiers (\texttt{nconst}), which are embedded in our person \acp{iri}. The \texttt{BIND(REPLACE(...))} pattern extracts the \texttt{nconst} at query time without materialising any \texttt{sameAsExternal} triples, demonstrating how Linked Data \acp{iri} designed as Cool~\acp{uri}~\cite{cooluris} can be re-used for federation across endpoints.

\begin{lstlisting}[language=SPARQL]
PREFIX ce:   <http://cineexplorer.local/ontology#>
PREFIX wdt:  <http://www.wikidata.org/prop/direct/>
PREFIX rdfs: <http://www.w3.org/2000/01/rdf-schema#>

SELECT ?localName ?nconst ?wikidataItem ?wikidataLabel ?birthDate
WHERE {
  VALUES ?person { <.../nm0000128> <.../nm0004266> ... } # 10 persons
  ?person ce:personName ?localName .
  BIND(REPLACE(STR(?person), "^.*/", "") AS ?nconst)
  SERVICE <https://query.wikidata.org/sparql> {
    ?wikidataItem wdt:P345 ?nconst .
    OPTIONAL { ?wikidataItem wdt:P569 ?birthDate . }
    OPTIONAL {
      ?wikidataItem rdfs:label ?wikidataLabel .
      FILTER(LANG(?wikidataLabel) = "en")
    }
  }
}
ORDER BY ?localName
\end{lstlisting}

\begin{table}[H]
\centering
\small
\begin{tabular}{llll}
\hline
\textbf{Local name} & \textbf{nconst} & \textbf{Wikidata} & \textbf{Born} \\
\hline
Alain Boublil           & nm0098842 & Q284540  & 1941-03-05 \\
Amanda Seyfried         & nm1086543 & Q189226  & 1985-12-03 \\
Anne Hathaway           & nm0004266 & Q36301   & 1982-11-12 \\
Claude-Michel Sch\"{o}nberg & nm0774744 & Q712004 & 1944-07-06 \\
Herbert Kretzmer        & nm0471014 & Q5734832 & 1925-10-05 \\
Hugh Jackman            & nm0413168 & Q129591  & 1968-10-12 \\
Russell Crowe           & nm0000128 & Q129817  & 1964-04-07 \\
Tom Hooper              & nm0393799 & Q295912  & 1972-10-05 \\
William Nicholson       & nm0629933 & Q706935  & 1948-01-12 \\
\hline
\end{tabular}
\caption{Q7 result: CineExplorer persons enriched with Wikidata birth dates (9 of 10 returned).}
\label{tab:q7-results}
\end{table}

Nine of the ten persons in the input \texttt{VALUES} list match (Table~\ref{tab:q7-results}). The tenth, Victor Hugo (\texttt{nm0401076}, Wikidata Q535), is correctly resolved on the Wikidata side --- but does not appear in our local \ac{kg} slice, so the join on \texttt{ce:personName} drops the row. This is the expected behaviour of a federated query: the local \ac{kg} anchors the join, and the SERVICE call enriches only those local resources that already exist. The match rate among persons that \emph{are} in the slice is therefore 100\%, reflecting Wikidata's comprehensive coverage of major film productions.

\subsection{Q8 — Seasons and Episodes per Series}

This query counts distinct seasons and episodes per series. The \texttt{OPTIONAL} clause handles episodes that lack a \texttt{ce:seasonNumber} value, producing \texttt{seasonCount~=~0} for those rows rather than excluding them.

\begin{lstlisting}[language=SPARQL]
PREFIX ce: <http://cineexplorer.local/ontology#>

SELECT ?seriesTitle
       (COUNT(DISTINCT ?season) AS ?seasonCount)
       (COUNT(DISTINCT ?ep)     AS ?episodeCount)
WHERE {
  ?series a ce:Series ;
          ce:workTitle ?seriesTitle ;
          ce:hasEpisode ?ep .
  OPTIONAL { ?ep ce:seasonNumber ?season . }
}
GROUP BY ?series ?seriesTitle
ORDER BY DESC(?episodeCount)
\end{lstlisting}

The query returns 16 series in the slice that have at least one linked episode. The migrated \ac{kg} contains 622 series and 139 episodes overall, but most series have only a small number of their episodes carried into the slice --- the slice is selected from \texttt{title.basics} by \texttt{numVotes} per title type, so only the most-voted-on individual episodes survive. The result is therefore dominated by a few cult series for which several episodes cleared the popularity threshold (e.g., \emph{Game of Thrones}, \emph{Breaking Bad}). The \texttt{OPTIONAL} clause cleanly returns \texttt{seasonCount~=~0} for episodes without a \texttt{ce:seasonNumber} value, demonstrating the gradual handling of missing properties under the \ac{owa}.

\subsection{Q9 — Directors Who Never Acted}

This query uses \texttt{MINUS} to subtract the set of directors who have an acting credit from the full set of directors.

\begin{lstlisting}[language=SPARQL]
PREFIX ce: <http://cineexplorer.local/ontology#>

SELECT ?name
WHERE {
  ?person a ce:Director ;
          ce:personName ?name .
  MINUS {
    ?person ce:actedIn ?anyTitle .
  }
}
ORDER BY ?name
\end{lstlisting}

\texttt{MINUS} operates on solution sets (whole result rows), while \texttt{FILTER NOT EXISTS} operates on variable bindings. For this pattern the two forms produce identical results, but \texttt{MINUS} expresses the intent more directly: remove from the director set those persons for whom an \texttt{actedIn} triple exists. The query returns 3{,}580 persons --- directors with no acting credit in the knowledge graph. The complementary set (directors who also acted) is 7{,}087~$-$~3{,}580~=~3{,}507, which agrees with the two-axis modelling of the \texttt{primaryProfession} attribute and the per-credit \texttt{Participation} role: a director may be typed as \texttt{ce:Director} either by Participation evidence (a \texttt{director} credit on some title) or by their \ac{imdb}-level career profession, and the same person can also accrue acting credits, which the query correctly excludes.

\subsection{Q10 — Titles Sharing Genre and Language}

Given \emph{Les Mis\'{e}rables} as target, this query finds titles that share at least one genre \emph{and} at least one language with it, ranked by the number of shared genres. An inner subquery identifies language-compatible candidates; the outer query then restricts to those sharing a genre and counts shared genres.

\begin{lstlisting}[language=SPARQL]
PREFIX ce: <http://cineexplorer.local/ontology#>

SELECT ?candidateTitle (COUNT(DISTINCT ?sharedGenre) AS ?commonGenres)
WHERE {
  BIND(<http://cineexplorer.local/data/title/tt1707386> AS ?target)
  ?target ce:hasGenre ?sharedGenre .
  ?candidate ce:hasGenre ?sharedGenre .
  FILTER(?candidate != ?target)
  ?candidate ce:workTitle ?candidateTitle .
  {
    SELECT DISTINCT ?candidate WHERE {
      BIND(<http://cineexplorer.local/data/title/tt1707386> AS ?target)
      ?target ce:language ?lang .
      ?candidate ce:language ?lang .
    }
  }
}
GROUP BY ?candidate ?candidateTitle
ORDER BY DESC(?commonGenres)
\end{lstlisting}

With \emph{Les Mis\'{e}rables} as target (3~genres, 7~languages), the query returns 2{,}751 candidates ranked by number of shared genres. The dense scale of the migrated \ac{kg} --- where many films share Drama as a genre and English or French as a language --- means a far larger candidate set than the curated sample produced. The subquery structure separates the language filter (inner) from the genre count (outer), demonstrating how subqueries can decompose multi-condition similarity patterns into reusable layers.
